%needs to be within 100 words
%To be suitable for Physical Review Letters (as opposed to a more specialized journal), the reported research should combine impact, innovation, and interest. When writing your paper, convey these elements clearly and directly. Impactful results are those that substantially advance a field, open a significant new area of research, or solve—or take a crucial step toward solving—a critical outstanding problem. Innovation reflects the extent to which a new result goes beyond previous work. And interest describes the intrinsic appeal of a result.

%At submission, authors should provide a 100-word compelling justification for why their paper meets Physical Review Letters’s criteria. This paragraph is a way for the authors to communicate directly and succinctly to the editors what they see as the main results, audience, and case for publication of the paper. It can aid in the choice of whether to send the paper out for review, and it can guide the choice of referees (particularly when the paper touches multiple fields).

%PRL Acceptance Criteria
%Submitted manuscripts should significantly advance fundamental or applied physical science by meeting one or more of the following criteria:

%Open a new research area, or a new avenue within an established area.
%Solve, or make essential steps towards solving, a critical problem.
%Introduce techniques or methods with highly significant impact.
%Be of unusual intrinsic interest to PRL's broad audience.

The Quantum Fourier Transform (QFT) is central to landmark quantum protocols like Shor's algorithm, but its $O(n^2)$ circuit depth present bottlenecks for near-term hardware. In this work, we identify that the QFT's ability to solve hidden subgroup problems relies on two fundamental properties: (i) maintaining measurements probability invariant under coset shifting and (ii) retaining discrete Fisher information about the target generator to ensure that the measurements contain learnable information about hidden subgroup. Guided by these requirements, we introduce HP-$1$ circuits with Hadamard and control phase gates that are friendly to current quantum devices.

Our numerical analysis proves that the HP-1 circuit maintain exponentially scaling Fisher information. After applying HP-1 circuit, we process the measurement outcomes with a polynomial-scaling neural network. We demonstrate our protocol on  Shor's algorithm, which successfully factoring the number $\leq 395$ with $11$ qubits circuit. The protocol shows robustness to noise in $9$ qubits cases. Thus, our protocol provides a robust, scalable pathway for executing advanced quantum algorithms on current hardware.

%GEMINI: Here is a revised significance statement tailored specifically for the PRL editors. It is exactly 98 words, strictly follows your formatting constraints, and directly addresses the PRL criteria of solving a critical problem and introducing a high-impact technique.

We introduce a new technique to solve an important problem facing quantum computation. The Quantum Fourier Transform (QFT) is used in Shor’s algorithm, but its O(n^2) depth limits near-term hardware implementation. We address this problem by distilling the QFT’s role in hidden subgroup problems, like Shor, to two fundamental properties: shift invariance and retention of Fisher information about the generator. Guided by this, we introduce hardware-friendly, O(n) HP-1 circuits. By combining these shallow circuits with efficient neural-network post-processing, we successfully numerically simulate Shor factoring with a shallow (depth 7) circuit (11 Hadamards and 30 control phase gates on 11 qubits).  